Volume 1 of the FCC Feasibility Report presents an overview of the physics case, experimental programme, and detector concepts for the Future Circular Collider (FCC). This volume outlines how  FCC would address some of the most profound open questions in particle physics, from precision studies of the Higgs and EW bosons and of the top quark, to the exploration of physics beyond the Standard Model. The report reviews the experimental opportunities offered by the staged implementation of  FCC, beginning with an electron-positron collider (FCC-ee), operating at several centre-of-mass energies, followed by a hadron collider (FCC-hh). Benchmark examples are given of the expected physics performance, in terms of precision and sensitivity to new phenomena, of each collider stage. Detector requirements and conceptual designs for FCC-ee experiments are discussed, as are the specific demands that the physics programme imposes on the accelerator in the domains of the calibration of the collision energy, and the interface region between the accelerator and the detector. The report also highlights advances in detector, software and computing technologies, as well as the theoretical tools /reconstruction techniques that will enable the precision measurements and discovery potential of the FCC experimental programme. The content and structure of this report are guided by the scope and priorities defined in the mandate of the FCC Feasibility Study. It is therefore not intended to serve as an exhaustive review of the full physics potential of FCC. Several topics, already covered in earlier reports such as the FCC CDR, are not reiterated here or are addressed only briefly, in alignment with the study’s focus. This volume reflects the outcome of a global collaborative effort involving hundreds of scientists and institutions, aided by a dedicated community-building coordination, and provides a targeted assessment of the scientific opportunities and experimental foundations of the FCC programme.